\begin{enumerate}
\item Los valores de las variables y su interpretación son los siguientes:

\begin{center}
\begin{tabular}{c|c|cccccccc|}
 & $z$ & $x_1$ & $x_2$ & $x_3$ & $x_4$ & $h_1$ & $h_3$ & $a_2$ & $a_3$\\ 
\hline
Phase 1 & 4 & $-1/5$ & 0 & $2/5$ & 0 & 0 & -1 & $-6/5$ & 0 \\ 
Phase 2 & -10000 & -250 & 0 & -100 & 0 & -100 & 0 & 0 & 0 \\ 
\hline
$x_4$ & 52  & $22/5$ & 0 & $16/5$ & 1 & 1 & 0 & $-3/5$ & 0\\ 
$x_2$ & 16  & $1/5$ & 1 & $3/5$ & 0 & 0 & 0 & $1/5$ & 0\\ 
$a_3$ & 4  & $-1/5$ & 0 & $2/5$ & 0 & 0 & -1 & $-1/5$ & 1\\ 
\hline
Phase 1 & 0 & 0 & 0 & 0 & 0 & 0 & 0 & -1 & -1 \\ 
Phase 2 & -9000 & -300 & 0 & 0 & 0 & -100 & -250 & -50 & 250 \\ 
\hline
$x_4$ & 20  & 6 & 0 & 0 & 1 & 1 & 8 & 1 & -8\\ 
$x_2$ & 10  & $1/2$ & 1 & 0 & 0 & 0 & $3/2$ & $1/2$ & $-3/2$\\ 
$x_3$ & 10  & $-1/2$ & 0 & 1 & 0 & 0 & $-5/2$ & $-1/2$ & $5/2$\\ 
\hline
\end{tabular}
\end{center}


\begin{itemize}
    \item $x_1 = 0$: no se produce aleación blanda. 
    \item $x_2 = 10$: se producen 10 toneladas de aleación dura.
    \item $x_3 = 10$: se producen 10 toneladas de aleación fuerte.
    \item $x_4 = 20$: se venden 20 toneladas de mineral rojo.
    \item $h_1 = 0$: todo el mineral rojo se emplea en aleaciones o bien se vende.
    \item $h_3 = 0$: se cumple el compromiso comercial por la cantidad mínima establecida (20 toneladas).
\end{itemize}

\item El precio sombra de la restricción 2 es $\pi^B_2 = 50$, con lo que por cada tonelada de mineral negro que la empresa pudiera extraer adicionalmente a las 80 que ya extrae el beneficio aumentaría en 50 euros, por lo que la empresa sí estaría interesada en extraer más que esas 80 toneladas, siempre y cuando el coste de extracción no superase un coste de 50 euros/tonelada.

\item Las columnas de las variables $h_1$ y $x_4$ en la matriz $A$ con idénticas: $A_{x_4} = A_{h_1} = (1\; 0\; 0)^T$ y sus respectivos criterios del simplex son:


$V^B_{x_4} = c_{x_4} - c^BB^{-1}A_{x_4} = 100 - c^BB^{-1}(1\; 0\; 0)^T$

$V^B_{h_1} = c_{h_1} - c^BB^{-1}A_{x_4} = 0 - c^BB^{-1}(1\; 0\; 0)^T$

Con lo que siempre se cumplirá $V^B_{x_4} > V^B_{h_1}$ y dado que las dos actividades son linealmente dependientes las dos no pueden ser simultáneamente básicas, así que la variable que será básica en la solución óptima será necesariamente $x_4$ y la variable $h_1$ será no básica y, por lo tanto, cero.

El términos de la situación de la empresa, lo anterior se pone de manifiesto en el hecho de que nunca le va a interesar dejar de utilizar mineral rojo ($h_1 > 0 $) porque la venta de ese mineral ($x_4 > 0$) permitirá mejorar el beneficio.

\item Para obtener el intervalo, se admite un valor del compromiso comercial genérico, $b_3$

\begin{equation}
\begin{split}
u^B = B^{-1}b = 
\begin{pmatrix}
1 & 1 & -8\\
0 & \frac{1}{2} & \frac{-3}{2}\\
0 & \frac{-1}{2} & \frac{5}{2}\\
\end{pmatrix}
\begin{pmatrix}
100\\
80\\
b_3\\
\end{pmatrix}
\geq 0 \Rightarrow
\begin{pmatrix}
180 - 8b_3\\
\frac{80-3b_3}{2}\\
\frac{-80+5b_3}{2}\\
\end{pmatrix}
\geq 0 \Rightarrow
16 \leq b_3 \leq \frac{45}{2} < \frac{80}{3}
\end{split}
\end{equation}

Mientras el compromiso comercial sea igual o mayor a 16 toneladas y menor o igual a 22.5 toneladas, las variables básicas de la solución son las misma y, por lo tanto, la gama de productos se mantiene.

\item $x_2 \leq 8 \Rightarrow x_2 + h_4 = 8$

\begin{center}
\begin{tabular}{c|c|ccccccccc|}
      & $z$   & $x_1$   & $x_2$   & $x_3$   & $x_4$ & $h_1$ & $h_3$ & $a_2$ & $a_3$ & $h_4$\\ 
      & -9000 & -300    & 0       & 0       & 0     & -100  & -250  & -50   & 250   & 0\\ 
\hline
$x_4$ & 20    & 6       & 0       & 0       & 1     & 1     & 8     & 1     & -8    & 0\\ 
$x_2$ & 10    & 1/2     & 1       & 0       & 0     & 0     & 3/2   & 1/2   & -3/2  & 0\\ 
$x_3$ & 10    & -1/2    & 0       & 1       & 0     & 0     & -5/2  & -1/2  & 5/2   & 0\\
\hline
      &  8    & 0       & 1       & 0       & 0     & 0     & 0     & 0     & 0     & 1\\
\hline
$h_4$ &  -2   & -1/2    & 0      & 0       & 0     & 0     & -3/2  & -1/2  & 3/2   & 1\\
\hline      
\end{tabular}
\end{center}

Podemos aplicar Lemke:

\begin{center}
\begin{tabular}{c|c|ccccccccc|}
      & $z$   & $x_1$   & $x_2$   & $x_3$   & $x_4$ & $h_1$ & $h_3$ & $a_2$ & $a_3$ & $h_4$\\ 
      & -9000 & -300    & 0       & 0       & 0     & -100  & -250  & -50   & 250   & 0\\ 
\hline
$x_4$ & 20    & 6       & 0       & 0       & 1     & 1     & 8     & 1     & -8    & 0\\ 
$x_2$ & 10    & 1/2     & 1       & 0       & 0     & 0     & 3/2   & 1/2   & -3/2  & 0\\ 
$x_3$ & 10    & -1/2    & 0       & 1       & 0     & 0     & -5/2  & -1/2  & 5/2   & 0\\
$h_4$ &  -2   & -1/2    & 0       & 0       & 0     & 0     & -3/2  & -1/2  & 3/2   & 1\\
\hline      
  &-26000/3   & -150    & -500/3  & 0       & 0     & 0     & 0     & 100/3 & 0     &  -500/3\\ 
\hline
$x_4$ &  28/3 &         &         &         &       &       & 0     &       &       & \\ 
$x_2$ &  8    &         &         &         &       &       & 0     &       &       &  \\ 
$x_3$ & 40/3  &         &         &         &       &       & 0     &       &       &  \\
$h_3$ & 4/3   & 1/3     & 0       & 0       & 0     & 0     & 1     & 1/3   & -1    & -2/3\\
\hline      
\end{tabular}
\end{center}

Los valores de las variables y su interpretación son los siguientes:
\begin{itemize}
    \item $x_1 = 0$: no se produce aleación blanda. 
    \item $x_2 = 5$: se producen 5 toneladas de aleación dura.
    \item $x_3 = 40/3$: se producen 40/3 toneladas de aleación fuerte.
    \item $x_4 = 28/3$: se venden 28/3 toneladas de mineral rojo.
    \item $h_1 = 0$: todo el mineral rojo se emplea en aleaciones o bien se vende.
    \item $h_3 = 4/3$: se cumple el compromiso comercial superando en 4/3 de tonelada la cantidad mínima establecida (20 toneladas).
\end{itemize}
    
Ahora el beneficio es menor e igual a 26000/3 euros.

\end{enumerate}